% \subsection{The Cartesian Product}
Given two sets $A$ and $B$, it is possible to "multiply" them to produce a new set denoted as $A \times B$. 
This operation is called the Cartesian product. To understand it, we must first understand the idea of an ordered pair. 
\begin{definition}
    [Ordered pair] An ordered pair is a list $\left(x, y\right)$ of two things $x$ and $y$, enclosed in parentheses and separated by a comma. 
\end{definition}
\begin{example}
    $(2, 4)$ is an ordered pair, as is $(4, 2)$. However, these two pairs are different, $(2, 4) \neq (4, 2)$.
\end{example}
Ordered pairs don't have to be numbersYou can have ordered pair of letters, sets, and matrix. 
\begin{example}
    Given $E = \left( \left\{2, 3\right\}, \left\{4, 5\right\}\right)$, $E$ is an ordered pair of sets. 
\end{example}

\begin{definition}
    [Cartesian product] The Cartesian product of two sets A and B is another set, donated as $A \times B$ and defined as $\left\{\left(a, b\right) :
    a \in A, b \in B\right\}$.
\end{definition}
For instance, given set $A = \left\{k, l\right\}$ and set $B = \left\{q, r\right\}$, the cartesian product of these two set is
$$A \times B = \left\{(k, q), (k, r), (l, q), (l, r)\right\}$$

Fun fact, if A and B are both finite sets, then 
\begin{equation*}
    \left|A \times B\right| = \left|A\right| \times \left|B\right|
\end{equation*}